\section{Contoh Program}
\label{sec:contoh-bab4}

Dua contoh berikut dipakai untuk menguji pemahaman tentang siklus
instruksi. Contoh pertama memakai program sesingkat mungkin agar
seluruh langkah dapat ditelusuri, sedangkan contoh kedua
membandingkan dua cara penjumlahan yang berbeda biayanya.

\subsection{Contoh 1: Menelusuri Siklus Instruksi sebuah Program}

Kode~\ref{lst:siklus-dasar} adalah program utuh yang berisi satu
instruksi utama, yaitu \lstasm|MOV AX, 1234h|. Program ini sengaja
dibuat pendek agar setiap langkah pada Emu8086 dapat diamati tanpa
kehilangan jejak.

\lstinputlisting[
    language=MASM,
    caption={Program \file{siklus-dasar.asm} untuk penelusuran siklus
    instruksi.},
    label={lst:siklus-dasar}
]{\CodeDir/siklus-dasar.asm}

Setelah program dirakit, hasil \eng{disassembly} pada Emu8086
memperlihatkan setiap instruksi beserta panjangnya dalam byte.
Panjang tersebut penting karena langsung menentukan berapa kali
\reg{IP} bertambah dan berapa banyak byte yang masuk ke
\eng{instruction queue}. Tabel~\ref{tab:trace-program} mencatat
penelusuran program tersebut dari titik masuk hingga program selesai.

\begin{table}[H]
    \centering
    \caption{Hasil penelusuran program \file{siklus-dasar.asm}
    dengan \eng{instruction queue} berisi maksimum enam byte.}
    \label{tab:trace-program}
    \small
    \begin{tabular}{@{}clp{1.5cm}p{2.3cm}p{2.0cm}p{3.6cm}@{}}
        \toprule
        \textbf{Langkah} & \textbf{\reg{CS}:\reg{IP}} &
        \textbf{Panjang} & \textbf{Instruksi} &
        \textbf{Antrean} & \textbf{Keadaan setelahnya} \\
        \midrule
        1 & \hex{1000}:\hex{0000} & 3 byte & \code{MOV AX, @DATA} &
        \hex{B8 00 10} & \reg{AX} $=$ \hex{1000} \\
        \addlinespace
        2 & \hex{1000}:\hex{0003} & 2 byte & \code{MOV DS, AX} &
        \hex{8E D8} & \reg{DS} $=$ \hex{1000} \\
        \addlinespace
        3 & \hex{1000}:\hex{0005} & 3 byte & \code{MOV AX, 1234h} &
        \hex{B8 34 12} & \reg{AX} $=$ \hex{1234} \\
        \addlinespace
        4 & \hex{1000}:\hex{0008} & 2 byte & \code{MOV AH, 4Ch} &
        \hex{B4 4C} & \reg{AH} $=$ \hex{4C} \\
        \addlinespace
        5 & \hex{1000}:\hex{000A} & 2 byte & \code{MOV AL, 00h} &
        \hex{B0 00} & \reg{AL} $=$ \hex{00} \\
        \addlinespace
        6 & \hex{1000}:\hex{000C} & 2 byte & \code{INT 21h} &
        \hex{CD 21} & Program berakhir \\
        \bottomrule
    \end{tabular}
\end{table}

Tiga hal dapat disimpulkan dari Tabel~\ref{tab:trace-program}.

\begin{enumerate}
    \item \textbf{\reg{IP} bertambah menurut panjang instruksi.}
          Instruksi tiga byte menambah \reg{IP} sebanyak tiga,
          sedangkan instruksi dua byte menambahnya sebanyak dua.

    \item \textbf{Antrean menampung lebih dari satu instruksi.}
          Pada langkah 4 dan 5, kedua instruksi hanya berukuran dua
          byte sehingga keduanya dapat berada dalam antrean enam
          byte sekaligus.

    \item \textbf{Instruksi yang mengakses memori menambah siklus
          bus.} Instruksi yang membaca variabel pada segmen data
          memicu siklus bus tambahan yang tidak dimiliki instruksi
          register-ke-register.
\end{enumerate}

\subsection{Contoh 2: Dua Cara Menjumlahkan Dua Bilangan}

Contoh kedua membandingkan biaya dua pendekatan penjumlahan. Pada
cara pertama, kedua operand sudah berada di register. Pada cara
kedua, operand dibaca dari memori sehingga BIU harus menjalankan
siklus bus. Program lengkapnya ditunjukkan pada
Kode~\ref{lst:jumlah-dua-cara}.

\lstinputlisting[
    language=MASM,
    caption={Program \file{jumlah-dua-cara.asm}: penjumlahan
    melalui register dan melalui memori.},
    label={lst:jumlah-dua-cara}
]{\CodeDir/jumlah-dua-cara.asm}

Kedua cara menghasilkan nilai yang sama, yaitu \hex{1234} $+$
\hex{5678} $=$ \hex{68AC}. Meskipun demikian, jumlah siklus yang
dibutuhkan berbeda. Tabel~\ref{tab:perbandingan-siklus} membandingkan
perkiraan jumlah siklus \eng{clock} untuk setiap instruksi
berdasarkan datasheet 8086 \cite{intel1979}.

\begin{table}[H]
    \centering
    \caption{Perkiraan jumlah siklus \eng{clock} kedua cara
    penjumlahan.}
    \label{tab:perbandingan-siklus}
    \small
    \begin{tabular}{@{}p{3.6cm}cp{4.4cm}p{3.4cm}@{}}
        \toprule
        \textbf{Instruksi} & \textbf{Siklus} & \textbf{Jenis operand} &
        \textbf{Akses bus} \\
        \midrule
        \multicolumn{4}{@{}l}{\textbf{Cara 1 --- operand di
        register}} \\
        \addlinespace
        \code{MOV AX, 1234h} & 4 & nilai segera &
        hanya pengambilan instruksi \\
        \addlinespace
        \code{MOV BX, 5678h} & 4 & nilai segera &
        hanya pengambilan instruksi \\
        \addlinespace
        \code{ADD AX, BX} & 3 & register--register &
        tidak ada \\
        \addlinespace
        \textbf{Jumlah cara 1} & \textbf{11} & & \\
        \midrule
        \multicolumn{4}{@{}l}{\textbf{Cara 2 --- operand di memori}} \\
        \addlinespace
        \code{MOV AX, bil1} & 14 & memori langsung &
        satu baca memori \\
        \addlinespace
        \code{ADD AX, bil2} & 15 & register--memori &
        satu baca memori \\
        \addlinespace
        \textbf{Jumlah cara 2} & \textbf{29} & & \\
        \bottomrule
    \end{tabular}
\end{table}

Perbandingan pada Tabel~\ref{tab:perbandingan-siklus} menunjukkan
bahwa cara kedua memerlukan waktu hampir tiga kali lebih lama untuk
pekerjaan yang sama. Penyebabnya bukan pada operasi penjumlahannya,
melainkan pada dua siklus bus tambahan yang harus dijalankan BIU.
Karena BIU juga sibuk mengambil instruksi berikutnya, siklus bus
yang panjang dapat mengosongkan antrean dan meniadakan keuntungan
\eng{pipeline}.

\begin{infobox}
Angka pada Tabel~\ref{tab:perbandingan-siklus} adalah nilai khas
untuk 8086 dengan pengalamatan langsung. Nilai sebenarnya bergantung
pada mode pengalamatan yang dipakai, karena setiap mode memiliki
biaya perhitungan alamat (\eng{effective address}) yang berbeda.
Cara penerjemahan instruksi menjadi kode mesin dibahas pada
Bab~\ref{chap:format-instruksi}.
\end{infobox}

\begin{session}[caption={Keluaran pemeriksaan isi memori
\file{jumlah-dua-cara.asm} pada Emu8086.}, label={lst:keluaran-hasil}]
1000:0008  hasil1  DW  68AC
1000:000A  hasil2  DW  68AC
\end{session}

Kedua penampung berisi nilai yang sama, tetapi waktu tempuhnya
berbeda. Inilah pelajaran utama bab ini: memahami arsitektur berarti
memahami \emph{mengapa} dua program yang hasilnya sama bisa memiliki
biaya yang berbeda.
